% ======================================================================
% Appendix A — h* observations
% ======================================================================

\appendix
\section{The \boldmath$h^*$-observations through \boldmath$n=10$}
\label{sec:hstar}

This appendix records the Ehrhart observations on the LR-zonotopes,
extended by this session to $n=9$ and $n=10$ in both core families. The
extension was executed as a one-day addendum with a narrow purpose:
establish whether the real-rootedness and log-concavity observed at
$n\le8$ continue, and whether the strictness decay and the
non-correlation with covering margins persist. Both answers are yes. The
appendix closes with the statement we regard as the residual question of
independent interest.

\subsection{The table}

Table~\ref{tab:hstar} lists, for each family and $n$: the
$h^*$-polynomial, real-rootedness (verified at $45$ decimal places on
the exact integer coefficients, with a companion-matrix check agreeing),
log-concavity and Newton's inequalities (exact rational arithmetic), and
the minimum ratio $h_k^2/(h_{k-1}h_{k+1})$, which is the
\emph{strictness} of log-concavity. Every row's Ehrhart data was
verified three ways (Eulerian expansion, standard inversion, and an
independent depth-first coset count at $t=1,2$), the volume identity
$e_{n-1}=\sum v_j$ was asserted exactly, and the loneliness identities
$\lambda=1/(n{+}1)$ (harmonic) and $\lambda=2/(2n{+}1)$ (rung-2) were
re-asserted exactly. The $n\le8$ rows reproduce the program's earlier
records digit for digit; the $n=9,10$ rows are new.

\begin{table}[htbp]
\centering
\footnotesize
\caption{$h^*$-polynomials of the LR-zonotopes, $n=3..10$, both core
families. ``RR'' = real-rooted (exact-integer verification at 45 dps);
``LC/Nw'' = log-concave and Newton (exact rational arithmetic);
$S$ = $\min_k h_k^2/(h_{k-1}h_{k+1})$; $nS$ = $n\cdot S$. The $n\le8$
rows reproduce the program's earlier records digit for digit; the
$n=9,10$ rows are new (this session).}
\label{tab:hstar}
\begin{tabularx}{\textwidth}{@{} l c >{\raggedright\arraybackslash}X c c r r @{}}
\toprule
family & $n$ & $h^*(t)$ & RR & LC/Nw & $S$ & $nS$ \\
\midrule
harm. & 3 & $1+7t+4t^2$ & y & y & $12.250$ & $36.8$ \\
rung2 & 3 & $1+11t+6t^2$ & y & y & $20.167$ & $60.5$ \\
harm. & 4 & $1+18t+35t^2+6t^3$ & y & y & $9.257$ & $37.0$ \\
rung2 & 4 & $1+22t+51t^2+10t^3$ & y & y & $9.490$ & $38.0$ \\
harm. & 5 & $1+37t+179t^2+133t^3+10t^4$ & y & y & $6.511$ & $32.6$ \\
rung2 & 5 & $1+45t+239t^2+181t^3+14t^4$ & y & y & $7.013$ & $35.1$ \\
harm. & 6 & $1+78t+744t^2+1286t^3+399t^4+12t^5$ & y & y & $5.518$
& $33.1$ \\
rung2 & 6 & $1+86t+920t^2+1682t^3+535t^4+16t^5$ & y & y & $5.748$
& $34.5$ \\
harm. & 7 & $1+147t+2522t^2+8948t^3+7323t^4$\newline
$+1201t^5+18t^6$ & y & y & $4.335$ & $30.3$ \\
rung2 & 7 & $1+161t+3024t^2+11144t^3+9295t^4$\newline
$+1551t^5+24t^6$ & y & y & $4.418$ & $30.9$ \\
harm. & 8 & $1+290t+8317t^2+51458t^3+82947t^4$\newline
$+35270t^5+3135t^6+22t^7$ & y & y & $3.791$ & $30.3$ \\
rung2 & 8 & $1+298t+9277t^2+60986t^3+102275t^4$\newline
$+44798t^5+4095t^6+30t^7$ & y & y & $3.829$ & $30.6$ \\
harm. & 9 & $1+559t+25519t^2+258901t^3+734059t^4$\newline
$+631993t^5+155305t^6+8035t^7+28t^8$ & y & y & $3.293$ & $29.6$ \\
rung2 & 9 & $1+585t+28885t^2+305035t^3+880149t^4$\newline
$+764143t^5+188659t^6+9789t^7+34t^8$ & y & y & $3.323$ & $29.9$ \\
harm. & 10 & $1+1098t+77550t^2+1213626t^3+5529648t^4$\newline
$+8309694t^5+4180482t^6+626910t^7+19359t^8+32t^9$ & y & y & $2.987$
& $29.9$ \\
rung2 & 10 & $1+1114t+83790t^2+1377794t^3+6457600t^4$\newline
$+9866694t^5+5017826t^6+758710t^7+23631t^8+40t^9$ & y & y & $3.004$
& $30.0$ \\
\bottomrule
\end{tabularx}
\end{table}

The sweep families behave identically: every instance of $(1,2,m)$,
$m\le14$, and $(1,2,3,m)$, $m\le12$, tested by the program was
real-rooted, log-concave, and Newton.

\subsection{Strictness decay and non-correlation}

Two quantitative facts were already in the program's records at $n\le8$
and persist through the new rows. First, the strictness $S$ decays
monotonically in $n$ in both families, with $n\cdot S$ drifting from
$\approx37$ at $n=3,4$ to $\approx29.6$--$30.0$ at $n=9,10$: the decay
is $\Theta(1/n)$ with a slowly drifting constant, so there is no growing
quantitative margin in the coefficient structure to spend. Second, the
strictness carries no correlation with the covering margins: across the
$n=3$ sweep (twelve instances, $(1,2,m)$, $m=3..14$) the Spearman
coefficient of strictness against the margin is $+0.32$; across the
$n=4$ sweep (eight instances, $(1,2,3,m)$, $m=5..12$) it is $-0.31$.
The sign flips between adjacent rungs; the harmonic family has margin
identically $0$ at every $n$ while its strictness varies by a factor of
four over the table. Both correlations were recomputed in the provenance
run from the persisted instance data and reproduce exactly.

The two facts close the bridge in the coefficient direction, and the
falsity of the covering-radius statement (\S\ref{sec:tm-zono}) closes it
in the metric direction: the statement that strict log-concavity would
have needed to feed is false at the extremal instances, and the
coefficient structure cannot see the metric hole depth that falsifies
it. There is no mechanism, and the arithmetic is clean enough that the
absence of the mechanism is itself a datum.

\subsection{The residual question}

What survives is a question of independent arithmetic interest, and we
state it decoupled from the Lonely Runner problem entirely. The
Ehrhart data of the LR-zonotopes is governed by the gcd structure
\eqref{eq:ehr} of the projected basis---the arithmetic matroid of the
configuration \cite{dAMoci}. In every instance tested (two core
families through $n=10$, sweeps through $m\le14$ at $n=3$ and $m\le12$
at $n=4$), the arithmetic $h^*$-polynomial is real-rooted. The natural
conjecture-shaped question is whether this holds for the projected
bases of the form $(1,\dots,n-1,m)$ or beyond, and the natural home for
it is the real-rootedness and Lorentzian literature for arithmetic
matroid Tutte specializations \cite{BrandenHuh,dAMoci,StanleyEC}. We
flag it as a question, record the data, and deliberately do not pursue
it further here: the extension to $n=9,10$ was executed to complete the
record, and the record's reading is that the question is real but
belongs to a different subject. It is not a bridge, and treating it as
one would repeat the pattern this paper documents.
